\section{HYM comparison with a divergence residual}\label{sec:comparison}

This section applies to any closed Riemann surface $C$ with a fixed
K\"ahler metric and any flat unitary bundle $(E,D,h_0)$ with parallel
determinant metric. A Hermitian metric is represented by a positive
$h_0$-self-adjoint endomorphism. The connection $D$, its formal adjoint
$D^*$, and the norms of endomorphism-valued forms always use this fixed
reference metric. Local parallel unitary frames have constant unitary
transition maps.

\subsection{The Chern equation and its trace inequality}
With area form $dA=a(z)\,dx\,dy=(i/2)a(z)\,dz\wedge d\bar z$, put
$\Delta=-4a^{-1}\partial_z\partial_{\bar z}$, the nonnegative
scalar Laplacian. The equation
\begin{equation}\label{eq:hym}
 h^{-1}h_s=-2i\Lambda F_h
\end{equation}
in a local parallel holomorphic frame is
\begin{equation}\label{eq:local-hym}
 h_s=4a^{-1}(h_{z\bar z}-h_{\bar z}h^{-1}h_z),
 \qquad
 2i\Lambda F_h=-4a^{-1}\partial_{\bar z}(h^{-1}h_z).
\end{equation}
In real coordinates it reads
\begin{equation}\label{eq:real-hym}
 \begin{split}
 ah_s={}&h_{xx}+h_{yy}-h_xh^{-1}h_x-h_yh^{-1}h_y\\
        &+i(h_xh^{-1}h_y-h_yh^{-1}h_x).
 \end{split}
\end{equation}
The last, Hermitian commutator term is retained throughout the proof.
For $h(x,y)=e^{xV+yW}$ at the origin, with noncommuting Hermitian
$V,W$, the first four terms cancel and this term is $i[V,W]$.
Thus \eqref{eq:hym} is not ordinary harmonic-map heat flow into
$\cP_r$.

For smooth positive initial data, the classical compact-base HYM
existence theorem gives a smooth solution for every finite time
\cite{donaldson}; a precise arbitrary-initial-metric statement
and continuation proof are in \cite[Theorem~6.8.25]{kobayashi}.
The book uses time $t=2s$ in our normalization; the statement
in \cite[\S4.1, following (4.8)]{hkkp} uses the normalization
above. The scalar equation for $\log\det h$
preserves the parallel determinant. The maximum principle applied to
$v^*hv$ for a constant local vector $v$ gives the upper order bound:
the quadratic term in \eqref{eq:local-hym} is nonpositive after this
pairing. For $h^{-1}$ the inverse derivative formula gives
\[
 (h^{-1})_s=4a^{-1}\bigl((h^{-1})_{z\bar z}
              -(h^{-1})_z h(h^{-1})_{\bar z}\bigr),
\]
which gives its upper bound as well. Hence
\begin{equation}\label{eq:order-bounds}
 ch_0\le h(0)\le C_0h_0
 \quad\Longrightarrow\quad
 ch_0\le h(s)\le C_0h_0\quad(s\ge0).
\end{equation}
These order comparisons are global and independent of spatial
derivatives of the initial metric. Existence is asserted separately
on each fixed compact smooth fibre; no uniform high-derivative
continuation theorem on degenerating fibres is needed.

Define the trace distance and the residual of a comparison metric by
\begin{equation}\label{eq:trace-residual}
 D_{\tr}(h,K)=\tr(h^{-1}K+K^{-1}h-2I),
 \qquad
 \cR(K)=K^{-1}K_s+2i\Lambda F_K.
\end{equation}
If the eigenvalues of $h^{-1}K$ are $e^{b_j}$, then
\begin{equation}\label{eq:distance-trace}
 d(h,K)^2=\sum_j b_j^2
       \le\sum_j(e^{b_j}+e^{-b_j}-2)=D_{\tr}(h,K).
\end{equation}
On a fixed bounded positive range the reverse inequality holds with
a fixed constant.

\begin{lemma}\label{lem:forced-trace}
For an exact solution $h$ and any smooth comparison metric $K$,
\begin{equation}\label{eq:forced-trace}
 (\partial_s+\Delta)D_{\tr}(h,K)
 \le\operatorname{Re}\tr\bigl((h^{-1}K-K^{-1}h)\cR(K)\bigr).
\end{equation}
\end{lemma}

\begin{proof}
First take two exact flows. At a specified point and time choose a
time-independent holomorphic frame normal for $h$, so $h=I$ and
$\partial h=0$ at that point. Substitution of the Chern equation
into the time derivative and scalar Laplacian of $\tr(h^{-1}K)$
cancels the curvature terms. The remaining expression is
$-4a^{-1}\tr(K_{\bar z}K^{-1}K_z)\le0$. Interchanging the two
metrics proves the inequality for the inverse trace. For a forced
comparison, its extra velocity is $K\cR(K)$. Its contributions to
these two traces are respectively
$\tr(h^{-1}K\cR(K))$ and $-\tr(\cR(K)K^{-1}h)$.
This proves \eqref{eq:forced-trace}. The real part is redundant for
the actual residual because $K\cR(K)$ is Hermitian; retaining it
allows complex pieces in a decomposition of that residual. The
calculation is invariant under changes of holomorphic frame and so
defines a global scalar inequality.
\end{proof}

\subsection{A comparison estimate using energy}
For an interval $[0,S]$ write
\[
 E_h(S)=\int_0^S\!\int_C|Dh|^2\,dA\,ds,
 \qquad E_K(S)=\int_0^S\!\int_C|DK|^2\,dA\,ds.
\]
The following lemma uses an actual global decomposition of the
residual. It does not assert that any small local error admits one.

\begin{lemma}\label{lem:weak-comparison}
Suppose $h,K$ have determinant equal to that of $h_0$ and satisfy
$cI\le h,K\le C_0I$. Suppose, in distributions on the closed
surface,
\begin{equation}\label{eq:divergence-residual}
 \cR(K)=D^*J+B,
\end{equation}
where $J$ is a complex endomorphism-valued one-form and $B$ an
endomorphism-valued function. Put
\[
 J_2(S)^2=\int_0^S\!\int_C|J|^2,\qquad
 B_1(S)=\int_0^S\!\int_C|B|,\qquad
 \eta=\int_C D_{\tr}(h(0),K(0)).
\]
For every $S>0$ for which these quantities are finite,
\begin{equation}\label{eq:weak-comparison}
 \sup_{0\le s\le S}\int_Cd(h(s),K(s))^2
 \le\eta+C B_1(S)
       +C\bigl(E_h(S)^{1/2}+E_K(S)^{1/2}\bigr)J_2(S).
\end{equation}
Moreover
\begin{equation}\label{eq:actual-energy}
 E_h(\infty)\le C\int_C\tr(h(0)-I)\,dA
                         \le C'\operatorname{Area}(C).
\end{equation}
The constants depend only on $c,C_0,r$. It is enough in
\eqref{eq:divergence-residual} that $J$ be piecewise $H^1$ with
zero normal traces at its internal boundaries, extended by zero
elsewhere. No spectral gap enters these estimates.
\end{lemma}

\begin{proof}
Put $Z=h^{-1}K-K^{-1}h$. On the bounded matrix range $|Z|\le C$.
The reference connection obeys the inverse and product rules, so
\begin{equation}\label{eq:dz}
 DZ=-h^{-1}(Dh)h^{-1}K+h^{-1}DK
           +K^{-1}(DK)K^{-1}h-K^{-1}Dh.
\end{equation}
In particular $|DZ|\le C(|Dh|+|DK|)$. Compatibility with trace
multiplication gives the bilinear integration-by-parts formula
\begin{equation}\label{eq:weak-pairing}
 \int_C\operatorname{Re}\tr(ZD^*J)
   =\int_C\operatorname{Re}\tr((DZ)\mathbin{\cdot}J).
\end{equation}
The dot contracts real cotangent indices; it does not reverse the
order of the matrices or insert an adjoint. Equivalently one can
apply the Hermitian formal-adjoint identity with test field $Z^*$.
For the stated piecewise fields, integration over each piece gives
the same formula because the normal traces vanish. Tangential jumps
do not contribute to divergence. More generally this is the defining
weak pairing whenever \eqref{eq:divergence-residual} holds with the
declared integrability.

Integrate \eqref{eq:forced-trace} on the closed surface and in
time. The scalar Laplacian integrates to zero. Use
\eqref{eq:weak-pairing}, \eqref{eq:dz} and Cauchy--Schwarz in
space-time for the two derivative terms. The $B$ term is at most
$CB_1(S)$. This first proves the bound for the trace distance;
\eqref{eq:distance-trace} yields \eqref{eq:weak-comparison}.

For the energy estimate, the exact equation \eqref{eq:local-hym}
gives
\begin{equation}\label{eq:trace-dissipation}
 -\frac{d}{ds}\int_C\tr(h-I)\,dA
       =4\int_C\tr(h_z^*h^{-1}h_z)\,dx\,dy.
\end{equation}
The derivative term integrates to zero; the displayed ordered
quadratic term is nonnegative. Only now replace $h^{-1}$ by its
scalar lower bound. The result controls a fixed multiple of
$\int_C|Dh|^2\,dA$. Determinant one implies $\tr(h-I)\ge0$
by the arithmetic--geometric mean inequality. Integrating to $S$
and then letting $S\to\infty$ proves \eqref{eq:actual-energy}.
It requires neither convergence to a particular parallel metric
nor a bound on the initial spatial derivative.
\end{proof}

In particular, on a bounded-area family, all-time approximation
follows if the initial trace error tends to zero, $E_K(\infty)$
stays bounded, and $J_2(\infty),B_1(\infty)$ tend to zero.
The remainder of the proof constructs exactly these quantities on
the compact family of Section~\ref{sec:geometry}.
